\documentclass[12pt]{article}
% Paquetes básicos primero
\usepackage[utf8]{inputenc}
\usepackage[T1]{fontenc}
\usepackage{ragged2e}
\usepackage{textcomp}
\usepackage{gensymb}
\usepackage{siunitx} 
\usepackage[spanish]{babel}
\usepackage{multicol}

% Geometría y diseño
\usepackage{geometry}
\geometry{
	letterpaper,
	left=2.5cm,
	right=2.5cm,
	top=3cm,
	bottom=3cm,
	headheight=60pt,
	headsep=1cm
}

% Paquetes de formato
\usepackage{setspace}
\usepackage{graphicx}
\usepackage{fancyhdr}
\usepackage{xcolor}
\usepackage{enumitem}
\usepackage{lastpage}
\usepackage{background}
\usepackage{listings}
\usepackage{tcolorbox}
\usepackage{fontawesome}
\usepackage{titlesec}
\usepackage{float}

% Paquetes para referencias
\usepackage[round]{natbib}  % Configuración específica para citas
\bibliographystyle{plainnat}  % Estilo de bibliografía

% URLs y enlaces (hyperref debe ir al final)
\usepackage{url}
\usepackage[hidelinks]{hyperref}

% Configuración de colores personalizados
\definecolor{tecnmBlue}{RGB}{0, 43, 89}
\definecolor{tecnmWhite}{RGB}{255, 255, 255}
\definecolor{codeBackground}{RGB}{240,240,240}
\definecolor{comentarios}{RGB}{63,127,95}
\definecolor{identificador}{RGB}{127,0,85}
\definecolor{cadena}{RGB}{42,0,255}
\definecolor{arduinoBlue}{RGB}{0,129,132}

% Configuración de listings para código Arduino
\lstdefinestyle{arduinoStyle}{
	language=C++,
	backgroundcolor=\color{codeBackground},
	basicstyle=\ttfamily\scriptsize,
	breaklines=true,
	captionpos=b,
	commentstyle=\color{comentarios},
	keywordstyle=\color{identificador},
	stringstyle=\color{cadena},
	numbers=left,
	numberstyle=\tiny\color{gray},
	stepnumber=1,
	numbersep=3pt,
	tabsize=3,
	showspaces=false,
	showstringspaces=false,
	frame=single,
	frameround=tttt,
	framesep=2mm,
	rulecolor=\color{arduinoBlue},
	emph={pinMode,digitalWrite,HIGH,LOW,INPUT,OUTPUT,tone,noTone,delay,delayMicroseconds,millis,INPUT_PULLUP},
	emphstyle={\color{blue}\bfseries},
	morekeywords={setup,loop},
	linewidth=\columnwidth,
	xleftmargin=0.5cm,
	xrightmargin=0.5cm,
	float=h,
	floatplacement=H,
	belowskip=-0.2cm,
	aboveskip=0.2cm
}

% Configuración de cajas personalizadas
\newtcolorbox{componentBox}[2][]{
	colback=tecnmBlue!5,
	colframe=tecnmBlue!40!black,
	title=#2,
	width=\columnwidth,
	#1
}

% Configuración del encabezado y pie de página
\pagestyle{fancy}
\fancyhf{}
\renewcommand{\headrulewidth}{0pt}
\renewcommand{\footrulewidth}{0pt}

% Definir comandos para reducir la carga
\newcommand{\smallgray}[1]{{\scriptsize\textcolor{gray}{#1}}}

\fancyhead[L]{%
	\hspace*{0.5cm}\raisebox{-2.2\height}{\includegraphics[height=1cm]{images/logotecnm.png}}%
}
\fancyhead[R]{%
	\raisebox{-1.7\height}{%
		\begin{tabular}[b]{r}
			{\scriptsize\textbf{Instituto Tecnológico de Morelia}}\\[-0.6em]
			\smallgray{Subdirección Académica}\\[-0.6em]
			\smallgray{Departamento de Sistemas y Computación}\\[-0.6em]
			\smallgray{Laboratorio de Sistemas Programables}\\[-0.4em]
			{\scriptsize Página \thepage\ de \pageref{LastPage}}
		\end{tabular}%
	}%
}

\makeatletter
\renewcommand\@biblabel[1]{#1.}
\makeatother

\fancyfoot[L]{%
	\small
	Implementación de un Buzzer
	
	\vspace{0.1cm}
	
	\small\textcolor{gray}{Plantilla de Reporte Laboratorio Sistemas Programables}
}
\fancyfoot[R]{%
	\small
	\textbf{Equipo 1}%
}

% Definir comando para citas
\newcommand{\cita}[1]{\cite{#1}}

% Estilo de títulos
\titleformat{\section}
{\color{tecnmBlue}\Large\bfseries}
{\thesection}{1em}{}

\titleformat{\subsection}
{\color{tecnmBlue}\large\bfseries}
{\thesubsection}{1em}{}

% Configuración de siunitx
\sisetup{
	detect-all,
	detect-family,
	per-mode=symbol,
	range-phrase=--,
	range-units=single,
	output-decimal-marker={,}
}

\begin{document}
	\pagenumbering{arabic}
	
	% Portada (una columna)
	% portada.tex
% Configuración del fondo solo para la portada
\backgroundsetup{
	scale=1,
	angle=0,
	opacity=0.2,
	contents={\includegraphics[width=\paperwidth,height=\paperheight]{images/fondo_portada.png}}
}
\BgThispage

\begin{center}
	\vspace*{2cm}
	
	\Large\textbf{Transistor BJT y MOSFET}
	
	\vspace{1.5cm}
	
	Autor(es):
	
	\vspace{0.5cm}
	
	Mario Eduardo Sánchez Mejía\\
	Fidel Alberto Zarco Áviles
	
	\vspace{0.3cm}
	
	\small{21120721@morelia.tecnm.mx}\\
	\small{l20121258@morelia.tecnm.mx}
	
	\vspace{1.5cm}
	
	Asesor(@s):
	
	\vspace{0.5cm}
	
	Luis Ulises Chávez Campos
	
	\vspace{0.5cm}
	
	\begin{center}
		\begin{abstract}
			\noindent
			\justifying
			Se presenta el desarrollo e implementación de un sistema de control electrónico basado en transistores como interruptores, utilizando microcontroladores Arduino. La implementación incorpora componentes específicos como el transistor BD235 y el MOSFET IRLZ14, junto con potenciómetros, para gestionar cargas de alta potencia, particularmente bombillas de 12V. El desarrollo experimental facilita la comprensión práctica del comportamiento de transistores en aplicaciones de conmutación, además de demostrar la implementación efectiva de sistemas de control mediante señales tanto digitales como analógicas. La integración resultante evidencia la aplicación práctica de principios fundamentales en electrónica de potencia y programación de microcontroladores en sistemas embebidos.
			
			\vspace{0.5cm}
			\noindent
			\textbf{Palabras clave:} Arduino, Transistores, MOSFET, Sistemas de Control, Electrónica de Potencia
		\end{abstract}
	\end{center}
\end{center}

% Desactivar el fondo para las siguientes páginas
\clearpage
\backgroundsetup{contents={}}
	% Inicio de dos columnas para el contenido principal
	\begin{multicols}{2}
		\setlength{\columnsep}{1cm} % Espacio entre columnas
		\tableofcontents
		\listoffigures
		\section{Introducción}
La presente documentación técnica establece los fundamentos teóricos y especificaciones técnicas de los componentes semiconductores y elementos pasivos utilizados en sistemas de control de cargas mediante transistores. El análisis abarca desde los principios fundamentales de la teoría de circuitos hasta las características específicas de los dispositivos semiconductores \citep{malvino2021electronic}.

\subsection{Transistor BJT BD235}
\begin{componentBox}{Especificaciones Técnicas y Configuración \citep{onsemi2021bd235}}
	\begin{itemize}[leftmargin=*,itemsep=1pt,parsep=1pt]
		\item \textbf{Clasificación}: Transistor Bipolar de Unión NPN de Potencia
		\item \textbf{Parámetros Característicos}:
		\begin{itemize}[itemsep=0pt,parsep=0pt]
			\item Tensión colector-emisor máxima (Vceo): \SI{50}{\volt}
			\item Corriente de colector máxima (Ic): \SI{2}{\ampere}
			\item Factor de amplificación de corriente (hFE): 40-250
			\item Configuración física: Encapsulado TO-126
		\end{itemize}
		\item \textbf{Configuración de Terminales}:
		\begin{itemize}[itemsep=0pt,parsep=0pt]
			\item Terminal de Colector (C): Electrodo principal de salida
			\item Terminal de Base (B): Electrodo de control
			\item Terminal de Emisor (E): Electrodo de referencia
		\end{itemize}
	\end{itemize}
\end{componentBox}

\subsection{MOSFET IRLZ44}
\begin{componentBox}{Características Técnicas y Operación \citep{infineon2020irlz44}}
	\begin{itemize}[leftmargin=*,itemsep=1pt,parsep=1pt]
		\item \textbf{Parámetros Operativos}:
		\begin{itemize}[itemsep=0pt,parsep=0pt]
			\item Tensión drenador-fuente máxima (Vds): \SI{55}{\volt}
			\item Corriente de drenador máxima (Id): \SI{47}{\ampere}
			\item Resistencia en conducción (Rds(on)): \SI{0.022}{\ohm}
			\item Tensión umbral de compuerta (Vgs(th)): \SI{1}{\volt}-\SI{2}{\volt}
		\end{itemize}
		\item \textbf{Estructura Terminal}:
		\begin{itemize}[itemsep=0pt,parsep=0pt]
			\item Terminal de Compuerta (G): Control de conducción
			\item Terminal de Drenador (D): Conducción principal
			\item Terminal de Fuente (S): Referencia de potencial
		\end{itemize}
	\end{itemize}
\end{componentBox}

\subsection{Potenciómetro de Precisión \SI{10}{\kilo\ohm}}
\begin{componentBox}{Parámetros Técnicos \citep{horowitz2015art}}
	\begin{itemize}[leftmargin=*,itemsep=1pt,parsep=1pt]
		\item \textbf{Especificaciones Fundamentales}:
		\begin{itemize}[itemsep=0pt,parsep=0pt]
			\item Resistencia nominal: \SI{10}{\kilo\ohm}
			\item Configuración mecánica: Sistema rotativo
			\item Característica de variación: Respuesta lineal
		\end{itemize}
		\item \textbf{Disposición de Terminales}:
		\begin{itemize}[itemsep=0pt,parsep=0pt]
			\item Terminal inicial: Punto de referencia
			\item Terminal variable: Cursor de control
			\item Terminal final: Límite resistivo
		\end{itemize}
	\end{itemize}
\end{componentBox}

\subsection{Sistema de Iluminación \SI{12}{\volt}}
\begin{componentBox}{Parámetros Operativos \citep{boylestad2013electronic}}
	\begin{itemize}[leftmargin=*,itemsep=1pt,parsep=1pt]
		\item \textbf{Especificaciones Eléctricas}:
		\begin{itemize}[itemsep=0pt,parsep=0pt]
			\item Tensión nominal de operación: \SI{12}{\volt}
			\item Potencia nominal: \SI{10}{\watt}
			\item Corriente de operación: \SI{0.83}{\ampere}
		\end{itemize}
	\end{itemize}
\end{componentBox}
		% explicaicon_codigo.tex

% Configuración preliminar para ajustar el espacio
\setlength{\parindent}{0pt}
\setlength{\parskip}{6pt}
\section{Explicación del Código}

\begin{figure}[H]
	\centering
	\begin{lstlisting}[
		language=C++,
		basicstyle=\ttfamily\scriptsize,
		numbers=left,
		numberstyle=\tiny\color{gray},
		backgroundcolor=\color{codeBackground},
		commentstyle=\color{comentarios},
		keywordstyle=\color{identificador},
		stringstyle=\color{cadena},
		emph={pinMode,digitalWrite,HIGH,LOW},
		emphstyle={\color{blue}\bfseries},
		frame=single,
		frameround=tttt,
		rulecolor=\color{arduinoBlue},
		breaklines=true,   % Permite la ruptura de líneas largas
		columns=fullflexible,
		keepspaces=true,
		basewidth=0.5em,
		xleftmargin=5mm,   % Márgenes laterales ajustados
		xrightmargin=0mm
		]
int tonePin = 4;
int ledMusic = 3;
unsigned long lastToneTime = 0;
bool isPlaying = false;
volatile bool shouldStopMelody = false;

void setup() {
	for (int pin = 5; pin <= 13; pin++) {
		pinMode(pin, INPUT);
	}
	pinMode(tonePin, OUTPUT);
	pinMode(ledMusic, OUTPUT);
}
	\end{lstlisting}
	\caption{Variables globales y función setup}
	\label{fig:codigo-led}
\end{figure}
La figura 2 muestra la configuración para el pin 4 para emitir sonidos y el pin 3 para controlar un LED indicador de música. En la función setup(), establece los pines 5 a 13 como entradas y define los pines tonePin y ledMusic como salidas para controlar el sonido y el LED. Además, usa variables para llevar el control del tiempo de reproducción y permite detener la melodía si es necesario.

		% explicaicon_codigo.tex

% Configuración preliminar para ajustar el espacio
\setlength{\parindent}{0pt}
\setlength{\parskip}{6pt}



\begin{figure}[H]
	\centering
	\begin{lstlisting}[
		language=C++,
		basicstyle=\ttfamily\scriptsize,
		numbers=left,
		numberstyle=\tiny\color{gray},
		backgroundcolor=\color{codeBackground},
		commentstyle=\color{comentarios},
		keywordstyle=\color{identificador},
		stringstyle=\color{cadena},
		emph={pinMode,digitalWrite,HIGH,LOW},
		emphstyle={\color{blue}\bfseries},
		frame=single,
		frameround=tttt,
		rulecolor=\color{arduinoBlue},
		breaklines=true,   % Permite la ruptura de líneas largas
		columns=fullflexible,
		keepspaces=true,
		basewidth=0.5em,
		xleftmargin=5mm,   % Márgenes laterales ajustados
		xrightmargin=0mm
		]
void loop() {
	bool anyKeyPressed = false;
	
	// Piano keys functionality
	if (digitalRead(13) == HIGH) {
		stopCurrentMelody();
		tone(tonePin, 262, 100);  // C4 (Do)
		anyKeyPressed = true;
		lastToneTime = millis();
	}
	if (digitalRead(12) == HIGH) {
		stopCurrentMelody();
		tone(tonePin, 294, 100);  //D4 (Re)
		anyKeyPressed = true;
		lastToneTime = millis();
	}
	\end{lstlisting}
	\caption{Función loop y condicionales}
	\label{fig:codigo-led}
\end{figure}
La figura 3  implementa la funcionalidad de un teclado de piano básico en el loop(). Define una variable anyKeyPressed para indicar si alguna tecla ha sido presionada.

Si el pin 13 detecta una señal alta (HIGH), el programa detiene la melodía actual (con stopCurrentMelody()), reproduce la nota C4 (Do) en el pin de tono (tonePin) a 262 Hz durante 100 ms, marca anyKeyPressed como verdadero, y actualiza el tiempo de la última nota tocada (lastToneTime). Del mismo modo, si el pin 12 está en HIGH, detiene la melodía, reproduce la nota D4 (Re) a 294 Hz durante 100 ms, y actualiza las mismas variables.

		% explicaicon_codigo.tex

% Configuración preliminar para ajustar el espacio
\setlength{\parindent}{0pt}
\setlength{\parskip}{6pt}
\begin{figure}[H]
	\centering
	\begin{lstlisting}[
		language=C++,
		basicstyle=\ttfamily\scriptsize,
		numbers=left,
		numberstyle=\tiny\color{gray},
		backgroundcolor=\color{codeBackground},
		commentstyle=\color{comentarios},
		keywordstyle=\color{identificador},
		stringstyle=\color{cadena},
		emph={pinMode,digitalWrite,HIGH,LOW},
		emphstyle={\color{blue}\bfseries},
		frame=single,
		frameround=tttt,
		rulecolor=\color{arduinoBlue},
		breaklines=true,   % Permite la ruptura de líneas largas
		columns=fullflexible,
		keepspaces=true,
		basewidth=0.5em,
		xleftmargin=5mm,   % Márgenes laterales ajustados
		xrightmargin=0mm
		]
if (digitalRead(11) == HIGH) {
	stopCurrentMelody();
	tone(tonePin, 330, 100);  // E4 (Mi)
	anyKeyPressed = true;
	lastToneTime = millis();
}
if (digitalRead(10) == HIGH) {
	stopCurrentMelody();
	tone(tonePin, 349, 100);  // F4 (Fa)
	anyKeyPressed = true;
	lastToneTime = millis();
}
if (digitalRead(9) == HIGH) {
	stopCurrentMelody();
	tone(tonePin, 392, 100);  // G4 (Sol)
	anyKeyPressed = true;
	lastToneTime = millis();
}
if (digitalRead(8) == HIGH) {
	stopCurrentMelody();
	tone(tonePin, 440, 100);  // A4 (La)
	anyKeyPressed = true;
	lastToneTime = millis();
}
if (digitalRead(7) == HIGH) {
	stopCurrentMelody();
	tone(tonePin, 494, 100);  // B4 (Si)
	anyKeyPressed = true;
	lastToneTime = millis();
}
	\end{lstlisting}
	\caption{Condicionales}
	\label{fig:codigo-led}
\end{figure}
La figura 4 amplía la funcionalidad del teclado de piano en el loop() de Arduino, permitiendo tocar más notas musicales.

Para cada pin del 11 al 7, el programa revisa si está en HIGH (tecla presionada). Si es así, detiene cualquier melodía en curso (stopCurrentMelody()), reproduce una nota específica (Mi, Fa, Sol, La o Si) en el pin de tono (tonePin) con una frecuencia de entre 330 y 494 Hz durante 100 ms, marca que una tecla fue presionada (anyKeyPressed = true), y actualiza el tiempo en lastToneTime.

		% explicaicon_codigo.tex

% Configuración preliminar para ajustar el espacio
\setlength{\parindent}{0pt}
\setlength{\parskip}{6pt}
\begin{figure}[H]
	\centering
	\begin{lstlisting}[
		language=C++,
		basicstyle=\ttfamily\scriptsize,
		numbers=left,
		numberstyle=\tiny\color{gray},
		backgroundcolor=\color{codeBackground},
		commentstyle=\color{comentarios},
		keywordstyle=\color{identificador},
		stringstyle=\color{cadena},
		emph={pinMode,digitalWrite,HIGH,LOW},
		emphstyle={\color{blue}\bfseries},
		frame=single,
		frameround=tttt,
		rulecolor=\color{arduinoBlue},
		breaklines=true,   % Permite la ruptura de líneas largas
		columns=fullflexible,
		keepspaces=true,
		basewidth=0.5em,
		xleftmargin=5mm,   % Márgenes laterales ajustados
		xrightmargin=0mm
		]
	if (digitalRead(6) == HIGH) {
		stopCurrentMelody();
		tone(tonePin, 523, 100);  // C5
		anyKeyPressed = true;
		lastToneTime = millis();
	}
	if (digitalRead(5) == HIGH) {
		shouldStopMelody = false;
		tetrisTheme();
	}
	
	\end{lstlisting}
	\caption{Condicionales}
	\label{fig:codigo-led}
\end{figure}
La figura 5 completa el teclado de piano y agrega una función para tocar una melodía específica.

Cuando el pin 6 está en HIGH, el programa detiene la melodía actual y reproduce la nota C5 (Do en una octava más alta) en el pin de tono (tonePin) a 523 Hz durante 100 ms. Marca que se ha presionado una tecla (anyKeyPressed = true) y actualiza el tiempo en lastToneTime.

Si el pin 5 está en HIGH, el programa habilita la reproducción de la melodía de Tetris (shouldStopMelody = false) y llama a la función tetrisTheme() para iniciar la melodía.
 
		% explicaicon_codigo.tex

% Configuración preliminar para ajustar el espacio
\setlength{\parindent}{0pt}
\setlength{\parskip}{6pt}
\begin{figure}[H]
	\centering
	\begin{lstlisting}[
		language=C++,
		basicstyle=\ttfamily\scriptsize,
		numbers=left,
		numberstyle=\tiny\color{gray},
		backgroundcolor=\color{codeBackground},
		commentstyle=\color{comentarios},
		keywordstyle=\color{identificador},
		stringstyle=\color{cadena},
		emph={pinMode,digitalWrite,HIGH,LOW},
		emphstyle={\color{blue}\bfseries},
		frame=single,
		frameround=tttt,
		rulecolor=\color{arduinoBlue},
		breaklines=true,   % Permite la ruptura de líneas largas
		columns=fullflexible,
		keepspaces=true,
		basewidth=0.5em,
		xleftmargin=5mm,   % Márgenes laterales ajustados
		xrightmargin=0mm
		]

	if (!anyKeyPressed && (millis() - lastToneTime > 100)) {
		digitalWrite(ledMusic, LOW);
		isPlaying = false;
	}
	
	delay(10);  
	}
	
		
	\end{lstlisting}
	\caption{Condicionales}
	\label{fig:codigo-led}
\end{figure}
La figura 6 muestra el apagado del LED indicador de música si no se está presionando ninguna tecla y han pasado al menos 100 ms desde la última nota tocada.

Primero, verifica si no hay teclas presionadas (!anyKeyPressed) y si el tiempo transcurrido desde la última nota (millis() - lastToneTime) supera los 100 ms. Si ambas condiciones se cumplen, apaga el LED (digitalWrite(ledMusic, LOW)) y marca que la música no está sonando (isPlaying = false). Finalmente, incluye un pequeño retraso de 10 ms para estabilizar el ciclo del loop().
		% explicaicon_codigo.tex

% Configuración preliminar para ajustar el espacio
\setlength{\parindent}{0pt}
\setlength{\parskip}{6pt}
\begin{figure}[H]
	\centering
	\begin{lstlisting}[
		language=C++,
		basicstyle=\ttfamily\scriptsize,
		numbers=left,
		numberstyle=\tiny\color{gray},
		backgroundcolor=\color{codeBackground},
		commentstyle=\color{comentarios},
		keywordstyle=\color{identificador},
		stringstyle=\color{cadena},
		emph={pinMode,digitalWrite,HIGH,LOW},
		emphstyle={\color{blue}\bfseries},
		frame=single,
		frameround=tttt,
		rulecolor=\color{arduinoBlue},
		breaklines=true,   % Permite la ruptura de líneas largas
		columns=fullflexible,
		keepspaces=true,
		basewidth=0.5em,
		xleftmargin=5mm,   % Márgenes laterales ajustados
		xrightmargin=0mm
		]
		
void stopCurrentMelody() {
	shouldStopMelody = true;
	noTone(tonePin);
	digitalWrite(ledMusic, LOW);
}

void toneLed(int frequency, int duration) {

	if (shouldStopMelody) {
		return;
	}
	

	for (int pin = 6; pin <= 13; pin++) {
		if (digitalRead(pin) == HIGH) {
			shouldStopMelody = true;
			return;
		}
	}
	
	digitalWrite(ledMusic, HIGH);
	tone(tonePin, frequency, duration);
	delay(duration);
	digitalWrite(ledMusic, LOW);
	delay(50);  
}


	
\end{lstlisting}
\caption{Funciones stopCurrentMelody() y toneLed()}
\label{fig:codigo-led}
\end{figure}
La figura 7 enseña la definición dos funciones en Arduino para controlar la reproducción de una melodía y un LED indicador. La función stopCurrentMelody() detiene la melodía actual, marcando la variable shouldStopMelody como true para indicar que la reproducción debe detenerse. Además, apaga el tono en el pin de sonido (tonePin) y el LED asociado (ledMusic).

La función toneLed(int frequency, int duration) reproduce un tono en el pin tonePin durante un tiempo específico. Si se ha indicado que la melodía debe detenerse (shouldStopMelody), la función termina sin hacer nada. Antes de iniciar el tono, también verifica si alguna tecla está presionada (pines del 6 al 13) y detiene la melodía si detecta que una tecla está activa. Al reproducir el tono, enciende el LED, mantiene el tono por el tiempo especificado, y luego apaga el LED. También incluye una breve pausa de 50 ms entre tonos para el parpadeo del LED.
		% explicaicon_codigo.tex

% Configuración preliminar para ajustar el espacio
\setlength{\parindent}{0pt}
\setlength{\parskip}{6pt}
\begin{figure}[H]
	\centering
	\begin{lstlisting}[
		language=C++,
		basicstyle=\ttfamily\scriptsize,
		numbers=left,
		numberstyle=\tiny\color{gray},
		backgroundcolor=\color{codeBackground},
		commentstyle=\color{comentarios},
		keywordstyle=\color{identificador},
		stringstyle=\color{cadena},
		emph={pinMode,digitalWrite,HIGH,LOW},
		emphstyle={\color{blue}\bfseries},
		frame=single,
		frameround=tttt,
		rulecolor=\color{arduinoBlue},
		breaklines=true,   % Permite la ruptura de líneas largas
		columns=fullflexible,
		keepspaces=true,
		basewidth=0.5em,
		xleftmargin=5mm,   % Márgenes laterales ajustados
		xrightmargin=0mm
		]
		
void tetrisTheme() {
	// Theme A
	toneLed(659, 250);  // E5
	if (shouldStopMelody) return;
	toneLed(494, 125);  // B4
	if (shouldStopMelody) return;
	toneLed(523, 125);  // C5
	if (shouldStopMelody) return;
	toneLed(587, 250);  // D5
	if (shouldStopMelody) return;
	toneLed(523, 125);  // C5
	if (shouldStopMelody) return;
	toneLed(494, 125);  // B4
	if (shouldStopMelody) return;
	toneLed(440, 250);  // A4
	if (shouldStopMelody) return;
	toneLed(440, 125);  // A4
	if (shouldStopMelody) return;
	toneLed(523, 125);  // C5
	if (shouldStopMelody) return;
	toneLed(659, 250);  // E5
	if (shouldStopMelody) return;
	toneLed(587, 125);  // D5
	if (shouldStopMelody) return;
	toneLed(523, 125);  // C5
	if (shouldStopMelody) return;
	toneLed(494, 375);  // B4
	if (shouldStopMelody) return;
	toneLed(523, 125);  // C5
	if (shouldStopMelody) return;
	toneLed(587, 250);  // D5
	if (shouldStopMelody) return;
	toneLed(659, 250);  // E5
	if (shouldStopMelody) return;
	toneLed(523, 250);  // C5
	if (shouldStopMelody) return;
	toneLed(440, 250);  // A4
	if (shouldStopMelody) return;
	toneLed(440, 500);  // A4
	if (shouldStopMelody) return;
	
	delay(250);  
	if (shouldStopMelody) return;
		
		
		
	\end{lstlisting}
	\caption{Función tetrisTheme() y condicionales}
	\label{fig:codigo-led}
\end{figure}
El código de la figura 8 define la función tetrisTheme() que reproduce el tema de Tetris (Theme A) utilizando tonos y frecuencias específicos. Cada llamada a toneLed() reproduce una nota musical en el pin designado (tonePin) con una duración específica, representando una parte de la melodía.

La función empieza reproduciendo una secuencia de notas con sus respectivas frecuencias y duraciones, como E5, B4, C5, y así sucesivamente. Después de cada nota, verifica si se ha activado shouldStopMelody, lo cual indica si la melodía debe detenerse y, de ser así, interrumpe la función.

La melodía incluye una breve pausa de 250 ms al final para permitir un intervalo entre repeticiones, haciendo que el tema se sienta más fluido y dándole tiempo al usuario antes de la próxima iteración.
		% explicaicon_codigo.tex

% Configuración preliminar para ajustar el espacio
\setlength{\parindent}{0pt}
\setlength{\parskip}{6pt}
\begin{figure}[H]
	\centering
	\begin{lstlisting}[
		language=C++,
		basicstyle=\ttfamily\scriptsize,
		numbers=left,
		numberstyle=\tiny\color{gray},
		backgroundcolor=\color{codeBackground},
		commentstyle=\color{comentarios},
		keywordstyle=\color{identificador},
		stringstyle=\color{cadena},
		emph={pinMode,digitalWrite,HIGH,LOW},
		emphstyle={\color{blue}\bfseries},
		frame=single,
		frameround=tttt,
		rulecolor=\color{arduinoBlue},
		breaklines=true,   % Permite la ruptura de líneas largas
		columns=fullflexible,
		keepspaces=true,
		basewidth=0.5em,
		xleftmargin=5mm,   % Márgenes laterales ajustados
		xrightmargin=0mm
		]
		
	// Theme B
	toneLed(494, 250);  // B4
	if (shouldStopMelody) return;
	toneLed(587, 250);  // D5
	if (shouldStopMelody) return;
	toneLed(659, 250);  // E5
	if (shouldStopMelody) return;
	toneLed(698, 250);  // F5
	if (shouldStopMelody) return;
	toneLed(659, 125);  // E5
	if (shouldStopMelody) return;
	toneLed(587, 125);  // D5
	if (shouldStopMelody) return;
	toneLed(523, 375);  // C5
	if (shouldStopMelody) return;
	toneLed(523, 125);  // C5
	if (shouldStopMelody) return;
	toneLed(587, 250);  // D5
	if (shouldStopMelody) return;
	toneLed(659, 250);  // E5
	if (shouldStopMelody) return;
	toneLed(523, 250);  // C5
	if (shouldStopMelody) return;
	toneLed(494, 250);  // B4
	if (shouldStopMelody) return;
	toneLed(440, 500);  // A4
	
	delay(500);  // Longer pause at the end
	}
\end{lstlisting}
\caption{Función tetrisTheme() y condicionales}
\label{fig:codigo-led}
\end{figure}
La figura 9 amplía la función tetrisTheme() en Arduino para incluir el "Theme B" de la melodía de Tetris. La función utiliza llamadas consecutivas a toneLed() para reproducir notas como B4, D5, E5, F5 y otras, cada una con una duración determinada. Después de cada nota, verifica la variable shouldStopMelody para ver si la reproducción debe detenerse; si es así, la función finaliza inmediatamente.

Al final de "Theme B", se agrega una pausa más larga de 500 ms para crear un descanso perceptible antes de cualquier repetición o cambio en la reproducción. Esto permite un flujo más natural entre los temas y evita que la melodía se sienta apresurada.
		%\input{sections/metodologia}
	\end{multicols}
	
	% Referencias (una columna)
	\clearpage
	\addcontentsline{toc}{section}{Referencias}
	\bibliography{referencias}
	
\end{document}